
\chapter{Methodik}

  %%%%%%%%%%%%%%%%%%%%
  % Achtung! Entwurf %
  %%%%%%%%%%%%%%%%%%%%

  Um die identifizierten Forschungslücken – insbesondere das Erklärungsdefizit und die mangelnde Nutzerzentrierung bestehender Systeme – systematisch zu adressieren, kombiniert diese Arbeit den \textit{Design Science Research (DSR)}-Ansatz \autocite{hevner2004} mit modernen Konzepten der Softwarearchitektur.
  Der Fokus liegt dabei auf drei Kernaspekten: der iterativen Prototypenentwicklung zur schrittweisen Verfeinerung der Anforderungen, der Implementierung einer modularen Systemarchitektur für maximale Anpassbarkeit, und der Integration quantitativer wie qualitativer Evaluationsmethoden.

  Nach einer detaillierten Darstellung des DSR-Rahmens mit seinen Designzyklen wird die technische Umsetzung erläutert, einschließlich der Auswahl geeigneter KI-Modelle, der Integration heterogener Datenquellen und der nutzerzentrierten Gestaltung der Erklärungsmechanismen.
  Den Abschluss bildet die Vorstellung eines mehrstufigen Evaluationskonzepts, das die Performance und weitere Metriken evaluiert.
  Diese strukturierte Vorgehensweise gewährleistet, dass das entwickelte Assistenzsystem nicht nur theoretisch fundiert, sondern auch praktisch robust und zukunftssicher ist.

\section{Design Science Research (DSR)-Ansatz}

    Der Design Science Research-Ansatz (DSR) verfolgt als forschungsleitendes Paradigma primär das Ziel, durch die Entwicklung neuartiger Artefakte praxisrelevante Probleme zu lösen.
    Im Unterschied zur verhaltenswissenschaftlichen Forschung, die auf Erklärungs- oder Prognosetheorien abzielt, steht bei DSR die Konstruktion nützlicher IT-Artefakte im Vordergrund, die die Grenzen menschlicher und organisatorischer Fähigkeiten erweitern können \autocite{hevner2004}.
    Hevner et al. betonen, dass im Design-Science-Paradigma Wissen über einen Problemraum gerade durch das Bauen und Anwenden des Artefakts gewonnen wird \autocite{hevner2004}.
    Ein DSR-Vorgehen ist somit ein prinzipieller Problemlösungsprozess, in dem \enquote{zweckgerichtete Artefakte} für spezifische Problem-domänen entwickelt werden müssen, um darin einen praktischen Nutzen zu stiften.
    Dabei sind Innovation und eine gründliche Evaluation des Artefakts zentrale Anforderungen \autocite{hevner2004}.

    Dieser Ansatz eignet sich besonders für die vorliegende Arbeit, da hier ein modularer Prototyp eines Finanzinformationssystems entwickelt wird, um identifizierte Forschungslücken im Bereich Robo-Advisors und MPT-gestützter Portfoliooptimierung zu schließen.
    Konkret wird zunächst eine Problemidentifikation durchgeführt, in der bestehende Systeme (Robo-Advisor, MPT-Tools) und die aus Kapitel 2 abgeleiteten Anforderungen analysiert werden.
    Anschließend folgt die Design- und Entwicklungsphase, in der ein artefaktspezifischer Prototyp unter iterativem Vorgehen erstellt wird.
    Dabei steuern Rückkopplungen mit Finanzexperten die sukzessive Verfeinerung des Moduls, was dem in der DSR-Literatur beschriebenen suchenden Entwicklungsprozess entspricht.
    Die abschließende Evaluationsphase validiert das Artefakt anhand definierter Nutzermetriken (Kapitel 5) und sichert so die praktische Relevanz und Wirksamkeit der Lösung.
    Insofern entspricht die methodische Vorgehensweise dem typischen DSR-Zyklus, bei dem Problemfindung, Konstruktion und Untersuchung des Artefakts eng verzahnt sind.

    Theoretisch wird das Design-Science-Vorgehen häufig als ein Zusammenspiel mehrerer Zyklen beschrieben.
    Hevner (2007) unterscheidet beispielsweise einen Relevanz-Zyklus, der Anforderungen aus dem Anwendungsumfeld in die Forschung einspeist, einen Rigor-Zyklus, der bewährte Theorien und Methoden einbindet, sowie einen mittleren Design-Zyklus, der die eigentliche Konstruktion und Evaluierung des Artefakts umfasst \autocite{hevner2007}.
    Der Relevanz-Zyklus stellt sicher, dass reale Problemstellungen berücksichtigt werden, der Rigor-Zyklus verankert das Vorgehen in der bestehenden Forschungsliteratur, und der zentrale Design-Zyklus widmet sich der konkreten Umsetzung und Prüfung der Lösung \autocite{hevner2007}.
    Diese Struktur verdeutlicht, dass DSR sowohl theoriegeleitet als auch anwendungsorientiert ist.

    Die in dieser Arbeit angewandten Phasen lassen sich an klassischen DSR-Modellen ablesen.
    Entsprechend den Vorschlägen von Peffers et al. werden die DSR-Aktivitäten typischerweise in die Schritte Problemidentifikation und -motivation, Zielformulierung, Konstruktion/Entwicklung, Demonstration (Nachweis der Nützlichkeit), Evaluation und Kommunikation unterteilt \autocite{peffers2007}.
    In der vorliegenden Arbeit wird der Schritt der Problemidentifikation durch die Analyse relevanter Finanzsysteme und identifizierter Forschungslücken umgesetzt.
    In der Konstruktionsphase entsteht der modulare Prototyp, der durch wiederholte Feedbackschleifen mit Experten ständig verfeinert wird.
    Schließlich werden im Demonstrations- und Evaluationsschritt die definierten Metriken (Kapitel 5) herangezogen, um die Leistungsfähigkeit und den praktischen Nutzen des Artefakts zu überprüfen.
    Somit gehen alle geplanten Arbeitsschritte auf bewährte DSR-Methoden zurück und sind iterativ angelegt.

    Wesentliche methodische Prinzipien des DSR sind dabei die enge Verknüpfung von Theorie und Praxis sowie die Fokussierung auf innovative Beiträge (Artefakte) und deren Qualität.
    Hevner et al. fordern, dass Design-Science-Artefakte sowohl neuartig (Guideline 4) als auch zweckgerichtet sein müssen und in ihrem Nutzen gründlich evaluiert werden \autocite{hevner2004}.
    Dies entspricht dem Vorgehen, wonach das entwickelte Finanzwerkzeug konkrete Verbesserungen gegenüber bestehenden Lösungen bieten soll und systematisch anhand festgelegter Kriterien bewertet wird.
    Der gesamte Forschungsprozess ist iterativ gestaltet: Erkenntnisse aus jeder Entwicklungs- und Testphase fließen in die Weiterentwicklung des Artefakts ein, was den konstruktiv-forschenden Charakter von DSR unterstreicht.
    Durch diese methodische Verankerung im DSR-Ansatz wird sichergestellt, dass die Arbeit nicht nur ein funktionierendes System erstellt, sondern auch wissenschaftlich relevante Erkenntnisse über den Lösungsprozess und das Design generiert \autocite{march1995}\autocite{hevner2004}.


\section{Modellauswahl und Modularität}

  \subsection{Entscheidungsmatrix für LLMs}

    Die Architektur folgt einem \textbf{modularen Ansatz}, der den Austausch von Modellen ohne Systemüberarbeitung ermöglicht.

    Grundlage bildete folgende Bewertung (nur Entwurf):
    \begin{table}[ht]
    \centering
    \caption{LLM-Auswahlkriterien für modulare Integration}
    \label{tab:llm-vergleich}
    \begin{tabular}{l|cccc}
      \textbf{Kriterium} & \textbf{DeepSeek R1} & \textbf{Llama 3} & \textbf{GPT-4o} & \textbf{Mixtral} \\
      \hline
      Open-Source & Ja & Ja & Nein & Ja \\
      API-Schnittstelle & REST & gRPC & REST & REST \\
      Finetuning-Support & PyTorch & HuggingFace & Nein & ONNX \\
      Token-Kosten & 0 & 0 & \$0.03/M & 0 \\
      \hline
      \textbf{Gesamt} & \textbf{1. Platz} & 2. Platz & 4. Platz & 3. Platz \\
    \end{tabular}
    \end{table}

  \subsection{Modulare Architektur}

    Das System implementiert drei Schlüsselkonzepte für Modellflexibilität:
    \begin{itemize}
      \item \textbf{Plugin-System}: Modelle werden als Pakete gekapselt (Standard-API-Schnittstelle)
      \item \textbf{Dynamic Routing}: Load-Balancer verteilt Anfragen basierend auf Modell-SLA (Latenz, Genauigkeit)
      \item \textbf{Ensemble Learning}: Kombination von Modellausgaben durch gewichtete Konsensalgorithmen \autocite{linardatos2021}
    \end{itemize}

    \begin{lstlisting}[language=Python,caption={Modell-Routing im Backend}]
    def get_ensemble_prediction(query):
        models = ["deepseek-r1", "mixtral", "llama3"]
        predictions = Parallel(n_jobs=3)(
            delayed(model_api.predict)(query, model)
            for model in models
        )
        return weighted_consensus(predictions)
    \end{lstlisting}

\section{Technischer Stack}

  \subsection{Backend (Flask)}

    \begin{itemize}
      \item \textbf{Modularität}: Flask-Blueprints für unabhängige Module (Datenpipeline, KI-Service, API)
      \item \textbf{API-Anbindung}: Möglichkeit größere und schnellere Online-Modelle (wie OpenAI oder Anbieter zu nutzen)
      \item \textbf{LLM-Orchestrierung}: Ollama-Integration für lokales Modellmanagement \autocite{wei2022}
      \item \textbf{Sicherheit}: JWT-Authentifizierung mit granularer Zugriffskontrolle
    \end{itemize}

  \subsection{Multi-Modell-Frontend}

    \begin{itemize}
      \item \textbf{Modellauswahl}: Nutzer können pro Anfrage Basismodell auswählen
      \item \textbf{Vergleichsmodus}: Side-by-Side-Darstellung verschiedener Modellantworten
      \item \textbf{Confidence Scoring}: Visualisierung von Unsicherheiten zwischen Modellen
    \end{itemize}

\section{Datenpipeline}

  \subsection{Quellen}

    \begin{itemize}
      \item \textbf{Strukturiert}: Yahoo Finance API (historische Kurse)
      \item \textbf{Unstrukturiert}: News und andere textuelle Informationen
    \end{itemize}

  % \subsection{Multi-Modell-Validierung}

  %   \begin{figure}[ht]
  %     \centering
  %     \includegraphics[width=0.9\textwidth]{data_validation_flow.pdf}
  %     \caption{Cross-Check von Modellaussagen durch unabhängige Datenquellen}
  %     \label{fig:validation_flow}
  %   \end{figure}

\section{Evaluationskriterien}

  \begin{table}[ht]
  \centering
  \caption{Modulspezifische Evaluationsmetriken}
  \begin{tabular}{lll}
    \textbf{Modul} & \textbf{KPI} & \textbf{Zielwert} \\
    \hline
    Modell-Interoperabilität & API-Latenz & tbd \\
    Multi-Modell-Konsens & Fleiss' Kappa & tbd \\
    Erklärbarkeit & SHAP-Wert-Konsistenz & tbd \\
    Skalierbarkeit & Requests/Min (50\% Load) & tbd \\
  \end{tabular}
  \end{table}

\section{Zusammenfassung}

  Die gewählte Methodik ermöglicht durch ihren modularen Aufbau:
  \begin{itemize}
    \item \textbf{Flexible Modellnutzung}: Austausch/Ergänzung von LLMs ohne Codeänderung
    \item \textbf{Plausibilitätschecks}: Cross-Validation durch parallele Modellausführungen
    \item \textbf{Zukunftssicherheit}: Einfache Integration kommender Modelle durch Flexibilität bei Modellauswahl
  \end{itemize}

%%%%%%%%%%%%%%%%%%%%%%%%%%%%%
\hspace{1.5cm}
\hrule


\begin{comment}
\textbf{To-Dos:}
\begin{itemize}
  \item Detaillierte Beschreibung des DSR-Ansatzes mit Design-Zyklen
  \item Nachvollziehbare Begründung der Modellauswahl (z. B, DeepSeek R1) und anderen technischen Gegebenheiten (z. B. Flask, Bootstrap etc)
  \item Nachvollziehbare Begründung des Prozesses, Aufbaus und Konzepts
  \begin{itemize}
    \item Welche Quellen werden gescrapt?
    \item Warum werden die Daten so ausgewertet wie sie sind?
  \end{itemize}
  \item Dokumentation der Evaluationskriterien vor Implementierung
  \item Klare Zuordnung von Methoden zu Forschungsfragen

  \item\textbf{Ideen:}
  \begin{itemize}
    \item Entscheidungsmatrix zur Framework-Auswahl (z. B. Ollama vs. Alternativen), Prozessabläufen etc
  \end{itemize}
\end{itemize}
\end{comment}